
\documentclass{article}
\usepackage{colortbl}
\usepackage{makecell}
\usepackage{multirow}
\usepackage{supertabular}

\begin{document}

\newcounter{utterance}

\centering \large Interaction Transcript for game `codenames', experiment `association\_difficult', episode 4 with deepseek{-}v4{-}pro.
\vspace{24pt}

{ \footnotesize  \setcounter{utterance}{1}
\setlength{\tabcolsep}{0pt}
\begin{supertabular}{c@{$\;$}|p{.15\linewidth}@{}p{.15\linewidth}p{.15\linewidth}p{.15\linewidth}p{.15\linewidth}p{.15\linewidth}}
    \# & $\;$A & \multicolumn{4}{c}{Game Master} & $\;\:$B\\
    \hline

    \theutterance \stepcounter{utterance}  
    & & & \multicolumn{2}{p{0.3\linewidth}}{
        \cellcolor[rgb]{0.9,0.9,0.9}{
            \makecell[{{p{\linewidth}}}]{
                \texttt{\tiny{[GM$|$GM]}}
                {'hidden': {'team': ['grožđe', 'stepenice', 'kaiš', 'temelj', 'pronevera', 'viola', 'gupi', 'puma', 'mandarina'], 'innocent': ['prevara', 'dukserica', 'pikolo', 'kajsija', 'koza', 'predsoblje', 'grouper'], 'opponent': ['pastrmka', 'mesec', 'šešir', 'tavan', 'rog', 'jagoda', 'magarac', 'milisekunda'], 'assassin': ['oposum']}, 'revealed': {'team': {'team': [], 'innocent': [], 'opponent': [], 'assassin': []}, 'opponent': {'team': [], 'innocent': [], 'opponent': [], 'assassin': []}}}
            }
        }
    }
    & & \\ \\

    \theutterance \stepcounter{utterance}  
    & & \multicolumn{4}{p{0.6\linewidth}}{
        \cellcolor[rgb]{0.9,0.9,0.9}{
            \makecell[{{p{\linewidth}}}]{
                \texttt{\tiny{[P1$\langle$GM]}}
                \texttt{Dajte jednu jednorečnu asocijaciju koja se odnosi na neke od reči vašeg tima sa liste ispod. Možete izabrati na koliko i na koje reči se vaša asocijacija odnosi.} \\
\texttt{Cilj je da vaš partner pogodi sve reči vašeg tima na osnovu asocijacija koje mu dajete, pre nego što drugi tim otkrije sve svoje. Ako vaš partner pogodi reč „protivnika" ili „distraktora", njegov potez pogađanja se završava i protivnički tim dobija priliku da odigra potez. Ako vaš partner pogodi reč „ubicu", igra se završava i vaš tim odmah gubi.} \\
\texttt{Što se vaša asocijacija odnosi na više reči, to brže možete pobediti, ali to partneru može biti teže da pogodi asocijacije, zato pažljivo birajte svoju reč{-}asocijaciju. Reč{-}asocijacija mora biti semantički povezana sa ciljanim rečima; ne može biti jedna od reči sa lista niti sadržati njihove delove.} \\
\\ 
\texttt{Uvek dajte svoju jednorečnu asocijaciju i svoju listu povezanih ciljanih reči razdvojenu zarezima u sledećem formatu i neka vaši odgovori budu što kraći; nikada ne uključujte nikakav drugi tekst osim onoga što je potrebno u ovom obliku:} \\
\\ 
\texttt{TRAG: <REČ>} \\
\texttt{CILJEVI: <CILJEVI>} \\
\\ 
\texttt{Reči vašeg tima su: grožđe, stepenice, kaiš, temelj, pronevera, viola, gupi, puma, mandarina.} \\
\texttt{Reči protivničkog tima su: pastrmka, mesec, šešir, tavan, rog, jagoda, magarac, milisekunda.} \\
\texttt{Reči{-}distraktori su: prevara, dukserica, pikolo, kajsija, koza, predsoblje, grouper.} \\
\texttt{Reči{-}ubice su: oposum.} \\
            }
        }
    }
    & \\ \\

    \theutterance \stepcounter{utterance}  
    & \multicolumn{4}{p{0.6\linewidth}}{
        \cellcolor[rgb]{0.9,0.9,0.9}{
            \makecell[{{p{\linewidth}}}]{
                \texttt{\tiny{[P1$\rangle$GM]}}
                \texttt{TRAG: voće} \\
\texttt{CILJEVI: grožđe, mandarina} \\
            }
        }
    }
    & & \\ \\

    \theutterance \stepcounter{utterance}  
    & & & \multicolumn{2}{p{0.3\linewidth}}{
        \cellcolor[rgb]{0.9,0.9,0.9}{
            \makecell[{{p{\linewidth}}}]{
                \texttt{\tiny{[GM$|$GM]}}
                \texttt{voće} \\
            }
        }
    }
    & & \\ \\

    \theutterance \stepcounter{utterance}  
    & & & \multicolumn{2}{p{0.3\linewidth}}{
        \cellcolor[rgb]{0.9,0.9,0.9}{
            \makecell[{{p{\linewidth}}}]{
                \texttt{\tiny{[GM$|$GM]}}
                ['grožđe', 'mandarina']
            }
        }
    }
    & & \\ \\

    \theutterance \stepcounter{utterance}  
    & & & \multicolumn{2}{p{0.3\linewidth}}{
        \cellcolor[rgb]{0.9,0.9,0.9}{
            \makecell[{{p{\linewidth}}}]{
                \texttt{\tiny{[GM$|$GM]}}
                {'word': 'grožđe', 'assignment': 'team'}
            }
        }
    }
    & & \\ \\

    \theutterance \stepcounter{utterance}  
    & & & \multicolumn{2}{p{0.3\linewidth}}{
        \cellcolor[rgb]{0.9,0.9,0.9}{
            \makecell[{{p{\linewidth}}}]{
                \texttt{\tiny{[GM$|$GM]}}
                {'word': 'mandarina', 'assignment': 'team'}
            }
        }
    }
    & & \\ \\

    \theutterance \stepcounter{utterance}  
    & & \multicolumn{4}{p{0.6\linewidth}}{
        \cellcolor[rgb]{0.9,0.9,0.9}{
            \makecell[{{p{\linewidth}}}]{
                \texttt{\tiny{[GM$\rangle$P2]}}
                \texttt{Navedite listu reči razdvojenih zarezima, do najviše 2 reči, iz sledeće liste koje se najbolje odnose ili su najbliže povezane sa rečju 'voće'. Uvek započnite svoju listu pogodaka sa POGODAK: i ne uključujte nikakav drugi tekst u svom odgovoru.} \\
\\ 
\texttt{viola, oposum, šešir, mesec, pikolo, temelj, stepenice, puma, kajsija, prevara, tavan, grouper, predsoblje, pronevera, koza, gupi, pastrmka, mandarina, jagoda, magarac, milisekunda, grožđe, kaiš, rog, dukserica} \\
            }
        }
    }
    & \\ \\

    \theutterance \stepcounter{utterance}  
    & & & \multicolumn{4}{p{0.6\linewidth}}{
        \cellcolor[rgb]{0.9,0.9,0.9}{
            \makecell[{{p{\linewidth}}}]{
                \texttt{\tiny{[GM$\langle$P2]}}
                \texttt{POGODAK: kajsija, jagoda} \\
            }
        }
    }
     \\ \\

    \theutterance \stepcounter{utterance}  
    & & & \multicolumn{2}{p{0.3\linewidth}}{
        \cellcolor[rgb]{0.9,0.9,0.9}{
            \makecell[{{p{\linewidth}}}]{
                \texttt{\tiny{[GM$|$GM]}}
                ['kajsija', 'jagoda']
            }
        }
    }
    & & \\ \\

    \theutterance \stepcounter{utterance}  
    & & & \multicolumn{2}{p{0.3\linewidth}}{
        \cellcolor[rgb]{0.9,0.9,0.9}{
            \makecell[{{p{\linewidth}}}]{
                \texttt{\tiny{[GM$|$GM]}}
                {'word': 'kajsija', 'assignment': 'innocent'}
            }
        }
    }
    & & \\ \\

    \theutterance \stepcounter{utterance}  
    & & & \multicolumn{2}{p{0.3\linewidth}}{
        \cellcolor[rgb]{0.9,0.9,0.9}{
            \makecell[{{p{\linewidth}}}]{
                \texttt{\tiny{[GM$|$GM]}}
                \texttt{kajsija} \\
            }
        }
    }
    & & \\ \\

    \theutterance \stepcounter{utterance}  
    & & & \multicolumn{2}{p{0.3\linewidth}}{
        \cellcolor[rgb]{0.9,0.9,0.9}{
            \makecell[{{p{\linewidth}}}]{
                \texttt{\tiny{[GM$|$GM]}}
                {'word': 'jagoda', 'assignment': 'opponent'}
            }
        }
    }
    & & \\ \\

    \theutterance \stepcounter{utterance}  
    & & & \multicolumn{2}{p{0.3\linewidth}}{
        \cellcolor[rgb]{0.9,0.9,0.9}{
            \makecell[{{p{\linewidth}}}]{
                \texttt{\tiny{[GM$|$GM]}}
                {'hidden': {'team': ['grožđe', 'stepenice', 'kaiš', 'temelj', 'pronevera', 'viola', 'gupi', 'puma', 'mandarina'], 'innocent': ['prevara', 'dukserica', 'pikolo', 'koza', 'predsoblje', 'grouper'], 'opponent': ['pastrmka', 'mesec', 'šešir', 'tavan', 'rog', 'magarac', 'milisekunda'], 'assassin': ['oposum']}, 'revealed': {'team': {'team': [], 'innocent': ['kajsija'], 'opponent': [], 'assassin': []}, 'opponent': {'team': [], 'innocent': [], 'opponent': ['jagoda'], 'assassin': []}}}
            }
        }
    }
    & & \\ \\

    \theutterance \stepcounter{utterance}  
    & & \multicolumn{4}{p{0.6\linewidth}}{
        \cellcolor[rgb]{0.9,0.9,0.9}{
            \makecell[{{p{\linewidth}}}]{
                \texttt{\tiny{[P1$\langle$GM]}}
                \texttt{Reč kajsija je pogođena, ali je to nevina reč. Potez tvog saigrača se tu završio.} \\
\texttt{Sada navedite još jedan trag koji se odnosi na neke od preostalih reči vašeg tima i listu povezanih ciljnih reči. Ne zaboravite da svoj trag započnete sa 'TRAG: ', zatim pređite u novi red i započnite svoju listu ciljnih reči, odvojenih zarezima, sa 'CILJEVI: '. Napomena: neke reči su uklonjene sa lista u poređenju sa prethodnim zahtevima.} \\
\\ 
\texttt{Vaše preostale reči tima su: grožđe, stepenice, kaiš, temelj, pronevera, viola, gupi, puma, mandarina.} \\
\texttt{Preostale reči vašeg protivnika su: pastrmka, mesec, šešir, tavan, rog, jagoda, magarac, milisekunda.} \\
\texttt{Preostale ometajuće reči su: prevara, dukserica, pikolo, koza, predsoblje, grouper.} \\
\texttt{Preostale reči ubice su: oposum.} \\
            }
        }
    }
    & \\ \\

    \theutterance \stepcounter{utterance}  
    & \multicolumn{4}{p{0.6\linewidth}}{
        \cellcolor[rgb]{0.9,0.9,0.9}{
            \makecell[{{p{\linewidth}}}]{
                \texttt{\tiny{[P1$\rangle$GM]}}
                \texttt{TRAG: riba} \\
\texttt{CILJEVI: gupi} \\
            }
        }
    }
    & & \\ \\

    \theutterance \stepcounter{utterance}  
    & & & \multicolumn{2}{p{0.3\linewidth}}{
        \cellcolor[rgb]{0.9,0.9,0.9}{
            \makecell[{{p{\linewidth}}}]{
                \texttt{\tiny{[GM$|$GM]}}
                \texttt{riba} \\
            }
        }
    }
    & & \\ \\

    \theutterance \stepcounter{utterance}  
    & & & \multicolumn{2}{p{0.3\linewidth}}{
        \cellcolor[rgb]{0.9,0.9,0.9}{
            \makecell[{{p{\linewidth}}}]{
                \texttt{\tiny{[GM$|$GM]}}
                ['gupi']
            }
        }
    }
    & & \\ \\

    \theutterance \stepcounter{utterance}  
    & & & \multicolumn{2}{p{0.3\linewidth}}{
        \cellcolor[rgb]{0.9,0.9,0.9}{
            \makecell[{{p{\linewidth}}}]{
                \texttt{\tiny{[GM$|$GM]}}
                {'word': 'gupi', 'assignment': 'team'}
            }
        }
    }
    & & \\ \\

    \theutterance \stepcounter{utterance}  
    & & \multicolumn{4}{p{0.6\linewidth}}{
        \cellcolor[rgb]{0.9,0.9,0.9}{
            \makecell[{{p{\linewidth}}}]{
                \texttt{\tiny{[GM$\rangle$P2]}}
                \texttt{Reč kajsija je pogođena, ali je to nevina reč. Tvoj potez se tu završio.} \\
\texttt{Sada navedite još jednu listu reči razdvojenih zarezima, od najmanje 1 do najviše 1 reči, iz sledeće liste reči koje se najbolje odnose ili su najuže povezane sa rečju 'riba'. Ne zaboravite da započnete svoj odgovor sa POGODAK:. Napomena: neke reči su uklonjene sa liste u odnosu na prethodne zahteve.} \\
\\ 
\texttt{viola, oposum, šešir, mesec, pikolo, temelj, stepenice, puma, prevara, tavan, grouper, predsoblje, pronevera, koza, gupi, pastrmka, mandarina, magarac, milisekunda, grožđe, kaiš, rog, dukserica} \\
            }
        }
    }
    & \\ \\

    \theutterance \stepcounter{utterance}  
    & & & \multicolumn{4}{p{0.6\linewidth}}{
        \cellcolor[rgb]{0.9,0.9,0.9}{
            \makecell[{{p{\linewidth}}}]{
                \texttt{\tiny{[GM$\langle$P2]}}
                \texttt{POGODAK: pastrmka} \\
            }
        }
    }
     \\ \\

    \theutterance \stepcounter{utterance}  
    & & & \multicolumn{2}{p{0.3\linewidth}}{
        \cellcolor[rgb]{0.9,0.9,0.9}{
            \makecell[{{p{\linewidth}}}]{
                \texttt{\tiny{[GM$|$GM]}}
                ['pastrmka']
            }
        }
    }
    & & \\ \\

    \theutterance \stepcounter{utterance}  
    & & & \multicolumn{2}{p{0.3\linewidth}}{
        \cellcolor[rgb]{0.9,0.9,0.9}{
            \makecell[{{p{\linewidth}}}]{
                \texttt{\tiny{[GM$|$GM]}}
                {'word': 'pastrmka', 'assignment': 'opponent'}
            }
        }
    }
    & & \\ \\

    \theutterance \stepcounter{utterance}  
    & & & \multicolumn{2}{p{0.3\linewidth}}{
        \cellcolor[rgb]{0.9,0.9,0.9}{
            \makecell[{{p{\linewidth}}}]{
                \texttt{\tiny{[GM$|$GM]}}
                \texttt{pastrmka} \\
            }
        }
    }
    & & \\ \\

    \theutterance \stepcounter{utterance}  
    & & & \multicolumn{2}{p{0.3\linewidth}}{
        \cellcolor[rgb]{0.9,0.9,0.9}{
            \makecell[{{p{\linewidth}}}]{
                \texttt{\tiny{[GM$|$GM]}}
                {'word': 'tavan', 'assignment': 'opponent'}
            }
        }
    }
    & & \\ \\

    \theutterance \stepcounter{utterance}  
    & & & \multicolumn{2}{p{0.3\linewidth}}{
        \cellcolor[rgb]{0.9,0.9,0.9}{
            \makecell[{{p{\linewidth}}}]{
                \texttt{\tiny{[GM$|$GM]}}
                {'hidden': {'team': ['grožđe', 'stepenice', 'kaiš', 'temelj', 'pronevera', 'viola', 'gupi', 'puma', 'mandarina'], 'innocent': ['prevara', 'dukserica', 'pikolo', 'koza', 'predsoblje', 'grouper'], 'opponent': ['mesec', 'šešir', 'rog', 'magarac', 'milisekunda'], 'assassin': ['oposum']}, 'revealed': {'team': {'team': [], 'innocent': ['kajsija'], 'opponent': ['pastrmka'], 'assassin': []}, 'opponent': {'team': [], 'innocent': [], 'opponent': ['jagoda', 'tavan'], 'assassin': []}}}
            }
        }
    }
    & & \\ \\

    \theutterance \stepcounter{utterance}  
    & & \multicolumn{4}{p{0.6\linewidth}}{
        \cellcolor[rgb]{0.9,0.9,0.9}{
            \makecell[{{p{\linewidth}}}]{
                \texttt{\tiny{[P1$\langle$GM]}}
                \texttt{Reč pastrmka je pogođena, ali je to protivnička reč. Potez tvog saigrača se tu završio.} \\
\texttt{Sada navedite još jedan trag koji se odnosi na neke od preostalih reči vašeg tima i listu povezanih ciljnih reči. Ne zaboravite da svoj trag započnete sa 'TRAG: ', zatim pređite u novi red i započnite svoju listu ciljnih reči, odvojenih zarezima, sa 'CILJEVI: '. Napomena: neke reči su uklonjene sa lista u poređenju sa prethodnim zahtevima.} \\
\\ 
\texttt{Vaše preostale reči tima su: grožđe, stepenice, kaiš, temelj, pronevera, viola, gupi, puma, mandarina.} \\
\texttt{Preostale reči vašeg protivnika su: mesec, šešir, tavan, rog, magarac, milisekunda.} \\
\texttt{Preostale ometajuće reči su: prevara, dukserica, pikolo, koza, predsoblje, grouper.} \\
\texttt{Preostale reči ubice su: oposum.} \\
            }
        }
    }
    & \\ \\

    \theutterance \stepcounter{utterance}  
    & \multicolumn{4}{p{0.6\linewidth}}{
        \cellcolor[rgb]{0.9,0.9,0.9}{
            \makecell[{{p{\linewidth}}}]{
                \texttt{\tiny{[P1$\rangle$GM]}}
                \texttt{TRAG: zloupotreba} \\
\texttt{CILJEVI: pronevera} \\
            }
        }
    }
    & & \\ \\

    \theutterance \stepcounter{utterance}  
    & & & \multicolumn{2}{p{0.3\linewidth}}{
        \cellcolor[rgb]{0.9,0.9,0.9}{
            \makecell[{{p{\linewidth}}}]{
                \texttt{\tiny{[GM$|$GM]}}
                \texttt{zloupotreba} \\
            }
        }
    }
    & & \\ \\

    \theutterance \stepcounter{utterance}  
    & & & \multicolumn{2}{p{0.3\linewidth}}{
        \cellcolor[rgb]{0.9,0.9,0.9}{
            \makecell[{{p{\linewidth}}}]{
                \texttt{\tiny{[GM$|$GM]}}
                ['pronevera']
            }
        }
    }
    & & \\ \\

    \theutterance \stepcounter{utterance}  
    & & & \multicolumn{2}{p{0.3\linewidth}}{
        \cellcolor[rgb]{0.9,0.9,0.9}{
            \makecell[{{p{\linewidth}}}]{
                \texttt{\tiny{[GM$|$GM]}}
                {'word': 'pronevera', 'assignment': 'team'}
            }
        }
    }
    & & \\ \\

    \theutterance \stepcounter{utterance}  
    & & \multicolumn{4}{p{0.6\linewidth}}{
        \cellcolor[rgb]{0.9,0.9,0.9}{
            \makecell[{{p{\linewidth}}}]{
                \texttt{\tiny{[GM$\rangle$P2]}}
                \texttt{Reč pastrmka je pogođena, ali je to protivnička reč. Tvoj potez se tu završio.} \\
\texttt{Sada navedite još jednu listu reči razdvojenih zarezima, od najmanje 1 do najviše 1 reči, iz sledeće liste reči koje se najbolje odnose ili su najuže povezane sa rečju 'zloupotreba'. Ne zaboravite da započnete svoj odgovor sa POGODAK:. Napomena: neke reči su uklonjene sa liste u odnosu na prethodne zahteve.} \\
\\ 
\texttt{viola, oposum, šešir, mesec, pikolo, temelj, stepenice, puma, prevara, grouper, predsoblje, pronevera, koza, gupi, mandarina, magarac, milisekunda, grožđe, kaiš, rog, dukserica} \\
            }
        }
    }
    & \\ \\

    \theutterance \stepcounter{utterance}  
    & & & \multicolumn{4}{p{0.6\linewidth}}{
        \cellcolor[rgb]{0.9,0.9,0.9}{
            \makecell[{{p{\linewidth}}}]{
                \texttt{\tiny{[GM$\langle$P2]}}
                \texttt{POGODAK: pronevera} \\
            }
        }
    }
     \\ \\

    \theutterance \stepcounter{utterance}  
    & & & \multicolumn{2}{p{0.3\linewidth}}{
        \cellcolor[rgb]{0.9,0.9,0.9}{
            \makecell[{{p{\linewidth}}}]{
                \texttt{\tiny{[GM$|$GM]}}
                ['pronevera']
            }
        }
    }
    & & \\ \\

    \theutterance \stepcounter{utterance}  
    & & & \multicolumn{2}{p{0.3\linewidth}}{
        \cellcolor[rgb]{0.9,0.9,0.9}{
            \makecell[{{p{\linewidth}}}]{
                \texttt{\tiny{[GM$|$GM]}}
                {'word': 'pronevera', 'assignment': 'team'}
            }
        }
    }
    & & \\ \\

    \theutterance \stepcounter{utterance}  
    & & & \multicolumn{2}{p{0.3\linewidth}}{
        \cellcolor[rgb]{0.9,0.9,0.9}{
            \makecell[{{p{\linewidth}}}]{
                \texttt{\tiny{[GM$|$GM]}}
                {'word': 'pronevera', 'assignment': 'team'}
            }
        }
    }
    & & \\ \\

    \theutterance \stepcounter{utterance}  
    & & & \multicolumn{2}{p{0.3\linewidth}}{
        \cellcolor[rgb]{0.9,0.9,0.9}{
            \makecell[{{p{\linewidth}}}]{
                \texttt{\tiny{[GM$|$GM]}}
                {'word': 'šešir', 'assignment': 'opponent'}
            }
        }
    }
    & & \\ \\

    \theutterance \stepcounter{utterance}  
    & & & \multicolumn{2}{p{0.3\linewidth}}{
        \cellcolor[rgb]{0.9,0.9,0.9}{
            \makecell[{{p{\linewidth}}}]{
                \texttt{\tiny{[GM$|$GM]}}
                {'hidden': {'team': ['grožđe', 'stepenice', 'kaiš', 'temelj', 'viola', 'gupi', 'puma', 'mandarina'], 'innocent': ['prevara', 'dukserica', 'pikolo', 'koza', 'predsoblje', 'grouper'], 'opponent': ['mesec', 'rog', 'magarac', 'milisekunda'], 'assassin': ['oposum']}, 'revealed': {'team': {'team': ['pronevera'], 'innocent': ['kajsija'], 'opponent': ['pastrmka'], 'assassin': []}, 'opponent': {'team': [], 'innocent': [], 'opponent': ['jagoda', 'tavan', 'šešir'], 'assassin': []}}}
            }
        }
    }
    & & \\ \\

    \theutterance \stepcounter{utterance}  
    & & \multicolumn{4}{p{0.6\linewidth}}{
        \cellcolor[rgb]{0.9,0.9,0.9}{
            \makecell[{{p{\linewidth}}}]{
                \texttt{\tiny{[P1$\langle$GM]}}
                \texttt{Reč pronevera je tačno pogođena. Potez tvog saigrača se tu završio.} \\
\texttt{Sada navedite još jedan trag koji se odnosi na neke od preostalih reči vašeg tima i listu povezanih ciljnih reči. Ne zaboravite da svoj trag započnete sa 'TRAG: ', zatim pređite u novi red i započnite svoju listu ciljnih reči, odvojenih zarezima, sa 'CILJEVI: '. Napomena: neke reči su uklonjene sa lista u poređenju sa prethodnim zahtevima.} \\
\\ 
\texttt{Vaše preostale reči tima su: grožđe, stepenice, kaiš, temelj, viola, gupi, puma, mandarina.} \\
\texttt{Preostale reči vašeg protivnika su: mesec, šešir, rog, magarac, milisekunda.} \\
\texttt{Preostale ometajuće reči su: prevara, dukserica, pikolo, koza, predsoblje, grouper.} \\
\texttt{Preostale reči ubice su: oposum.} \\
            }
        }
    }
    & \\ \\

    \theutterance \stepcounter{utterance}  
    & \multicolumn{4}{p{0.6\linewidth}}{
        \cellcolor[rgb]{0.9,0.9,0.9}{
            \makecell[{{p{\linewidth}}}]{
                \texttt{\tiny{[P1$\rangle$GM]}}
                \texttt{TRAG: instrument} \\
\texttt{CILJEVI: viola, pikolo} \\
            }
        }
    }
    & & \\ \\

    \theutterance \stepcounter{utterance}  
    & & & \multicolumn{2}{p{0.3\linewidth}}{
        \cellcolor[rgb]{0.9,0.9,0.9}{
            \makecell[{{p{\linewidth}}}]{
                \texttt{\tiny{[GM$|$GM]}}
                \texttt{instrument} \\
            }
        }
    }
    & & \\ \\

    \theutterance \stepcounter{utterance}  
    & & & \multicolumn{2}{p{0.3\linewidth}}{
        \cellcolor[rgb]{0.9,0.9,0.9}{
            \makecell[{{p{\linewidth}}}]{
                \texttt{\tiny{[GM$|$GM]}}
                ['viola', 'pikolo']
            }
        }
    }
    & & \\ \\

    \theutterance \stepcounter{utterance}  
    & & & \multicolumn{2}{p{0.3\linewidth}}{
        \cellcolor[rgb]{0.9,0.9,0.9}{
            \makecell[{{p{\linewidth}}}]{
                \texttt{\tiny{[GM$|$GM]}}
                {'word': 'viola', 'assignment': 'team'}
            }
        }
    }
    & & \\ \\

    \theutterance \stepcounter{utterance}  
    & & & \multicolumn{2}{p{0.3\linewidth}}{
        \cellcolor[rgb]{0.9,0.9,0.9}{
            \makecell[{{p{\linewidth}}}]{
                \texttt{\tiny{[GM$|$GM]}}
                {'word': 'pikolo', 'assignment': 'innocent'}
            }
        }
    }
    & & \\ \\

    \theutterance \stepcounter{utterance}  
    & & \multicolumn{4}{p{0.6\linewidth}}{
        \cellcolor[rgb]{0.9,0.9,0.9}{
            \makecell[{{p{\linewidth}}}]{
                \texttt{\tiny{[GM$\rangle$P2]}}
                \texttt{Reč pronevera je tačno pogođena. Tvoj potez se tu završio.} \\
\texttt{Sada navedite još jednu listu reči razdvojenih zarezima, od najmanje 1 do najviše 2 reči, iz sledeće liste reči koje se najbolje odnose ili su najuže povezane sa rečju 'instrument'. Ne zaboravite da započnete svoj odgovor sa POGODAK:. Napomena: neke reči su uklonjene sa liste u odnosu na prethodne zahteve.} \\
\\ 
\texttt{viola, oposum, mesec, pikolo, temelj, stepenice, puma, prevara, grouper, predsoblje, koza, gupi, mandarina, magarac, milisekunda, grožđe, kaiš, rog, dukserica} \\
            }
        }
    }
    & \\ \\

    \theutterance \stepcounter{utterance}  
    & & & \multicolumn{4}{p{0.6\linewidth}}{
        \cellcolor[rgb]{0.9,0.9,0.9}{
            \makecell[{{p{\linewidth}}}]{
                \texttt{\tiny{[GM$\langle$P2]}}
                \texttt{POGODAK: viola, pikolo} \\
            }
        }
    }
     \\ \\

    \theutterance \stepcounter{utterance}  
    & & & \multicolumn{2}{p{0.3\linewidth}}{
        \cellcolor[rgb]{0.9,0.9,0.9}{
            \makecell[{{p{\linewidth}}}]{
                \texttt{\tiny{[GM$|$GM]}}
                ['viola', 'pikolo']
            }
        }
    }
    & & \\ \\

    \theutterance \stepcounter{utterance}  
    & & & \multicolumn{2}{p{0.3\linewidth}}{
        \cellcolor[rgb]{0.9,0.9,0.9}{
            \makecell[{{p{\linewidth}}}]{
                \texttt{\tiny{[GM$|$GM]}}
                {'word': 'viola', 'assignment': 'team'}
            }
        }
    }
    & & \\ \\

    \theutterance \stepcounter{utterance}  
    & & & \multicolumn{2}{p{0.3\linewidth}}{
        \cellcolor[rgb]{0.9,0.9,0.9}{
            \makecell[{{p{\linewidth}}}]{
                \texttt{\tiny{[GM$|$GM]}}
                {'word': 'viola', 'assignment': 'team'}
            }
        }
    }
    & & \\ \\

    \theutterance \stepcounter{utterance}  
    & & & \multicolumn{2}{p{0.3\linewidth}}{
        \cellcolor[rgb]{0.9,0.9,0.9}{
            \makecell[{{p{\linewidth}}}]{
                \texttt{\tiny{[GM$|$GM]}}
                {'word': 'pikolo', 'assignment': 'innocent'}
            }
        }
    }
    & & \\ \\

    \theutterance \stepcounter{utterance}  
    & & & \multicolumn{2}{p{0.3\linewidth}}{
        \cellcolor[rgb]{0.9,0.9,0.9}{
            \makecell[{{p{\linewidth}}}]{
                \texttt{\tiny{[GM$|$GM]}}
                {'word': 'pikolo', 'assignment': 'innocent'}
            }
        }
    }
    & & \\ \\

    \theutterance \stepcounter{utterance}  
    & & & \multicolumn{2}{p{0.3\linewidth}}{
        \cellcolor[rgb]{0.9,0.9,0.9}{
            \makecell[{{p{\linewidth}}}]{
                \texttt{\tiny{[GM$|$GM]}}
                \texttt{pikolo} \\
            }
        }
    }
    & & \\ \\

    \theutterance \stepcounter{utterance}  
    & & & \multicolumn{2}{p{0.3\linewidth}}{
        \cellcolor[rgb]{0.9,0.9,0.9}{
            \makecell[{{p{\linewidth}}}]{
                \texttt{\tiny{[GM$|$GM]}}
                {'word': 'magarac', 'assignment': 'opponent'}
            }
        }
    }
    & & \\ \\

    \theutterance \stepcounter{utterance}  
    & & & \multicolumn{2}{p{0.3\linewidth}}{
        \cellcolor[rgb]{0.9,0.9,0.9}{
            \makecell[{{p{\linewidth}}}]{
                \texttt{\tiny{[GM$|$GM]}}
                {'hidden': {'team': ['grožđe', 'stepenice', 'kaiš', 'temelj', 'gupi', 'puma', 'mandarina'], 'innocent': ['prevara', 'dukserica', 'koza', 'predsoblje', 'grouper'], 'opponent': ['mesec', 'rog', 'milisekunda'], 'assassin': ['oposum']}, 'revealed': {'team': {'team': ['pronevera', 'viola'], 'innocent': ['kajsija', 'pikolo'], 'opponent': ['pastrmka'], 'assassin': []}, 'opponent': {'team': [], 'innocent': [], 'opponent': ['jagoda', 'tavan', 'šešir', 'magarac'], 'assassin': []}}}
            }
        }
    }
    & & \\ \\

    \theutterance \stepcounter{utterance}  
    & & \multicolumn{4}{p{0.6\linewidth}}{
        \cellcolor[rgb]{0.9,0.9,0.9}{
            \makecell[{{p{\linewidth}}}]{
                \texttt{\tiny{[P1$\langle$GM]}}
                \texttt{Reč viola je tačno pogođena. Reč pikolo je pogođena, ali je to nevina reč. Potez tvog saigrača se tu završio.} \\
\texttt{Sada navedite još jedan trag koji se odnosi na neke od preostalih reči vašeg tima i listu povezanih ciljnih reči. Ne zaboravite da svoj trag započnete sa 'TRAG: ', zatim pređite u novi red i započnite svoju listu ciljnih reči, odvojenih zarezima, sa 'CILJEVI: '. Napomena: neke reči su uklonjene sa lista u poređenju sa prethodnim zahtevima.} \\
\\ 
\texttt{Vaše preostale reči tima su: grožđe, stepenice, kaiš, temelj, gupi, puma, mandarina.} \\
\texttt{Preostale reči vašeg protivnika su: mesec, rog, magarac, milisekunda.} \\
\texttt{Preostale ometajuće reči su: prevara, dukserica, koza, predsoblje, grouper.} \\
\texttt{Preostale reči ubice su: oposum.} \\
            }
        }
    }
    & \\ \\

    \theutterance \stepcounter{utterance}  
    & \multicolumn{4}{p{0.6\linewidth}}{
        \cellcolor[rgb]{0.9,0.9,0.9}{
            \makecell[{{p{\linewidth}}}]{
                \texttt{\tiny{[P1$\rangle$GM]}}
                \texttt{TRAG: penjanje} \\
\texttt{CILJEVI: stepenice, puma} \\
            }
        }
    }
    & & \\ \\

    \theutterance \stepcounter{utterance}  
    & & & \multicolumn{2}{p{0.3\linewidth}}{
        \cellcolor[rgb]{0.9,0.9,0.9}{
            \makecell[{{p{\linewidth}}}]{
                \texttt{\tiny{[GM$|$GM]}}
                \texttt{penjanje} \\
            }
        }
    }
    & & \\ \\

    \theutterance \stepcounter{utterance}  
    & & & \multicolumn{2}{p{0.3\linewidth}}{
        \cellcolor[rgb]{0.9,0.9,0.9}{
            \makecell[{{p{\linewidth}}}]{
                \texttt{\tiny{[GM$|$GM]}}
                ['stepenice', 'puma']
            }
        }
    }
    & & \\ \\

    \theutterance \stepcounter{utterance}  
    & & & \multicolumn{2}{p{0.3\linewidth}}{
        \cellcolor[rgb]{0.9,0.9,0.9}{
            \makecell[{{p{\linewidth}}}]{
                \texttt{\tiny{[GM$|$GM]}}
                {'word': 'stepenice', 'assignment': 'team'}
            }
        }
    }
    & & \\ \\

    \theutterance \stepcounter{utterance}  
    & & & \multicolumn{2}{p{0.3\linewidth}}{
        \cellcolor[rgb]{0.9,0.9,0.9}{
            \makecell[{{p{\linewidth}}}]{
                \texttt{\tiny{[GM$|$GM]}}
                {'word': 'puma', 'assignment': 'team'}
            }
        }
    }
    & & \\ \\

    \theutterance \stepcounter{utterance}  
    & & \multicolumn{4}{p{0.6\linewidth}}{
        \cellcolor[rgb]{0.9,0.9,0.9}{
            \makecell[{{p{\linewidth}}}]{
                \texttt{\tiny{[GM$\rangle$P2]}}
                \texttt{Reč viola je tačno pogođena. Reč pikolo je pogođena, ali je to nevina reč. Tvoj potez se tu završio.} \\
\texttt{Sada navedite još jednu listu reči razdvojenih zarezima, od najmanje 1 do najviše 2 reči, iz sledeće liste reči koje se najbolje odnose ili su najuže povezane sa rečju 'penjanje'. Ne zaboravite da započnete svoj odgovor sa POGODAK:. Napomena: neke reči su uklonjene sa liste u odnosu na prethodne zahteve.} \\
\\ 
\texttt{oposum, mesec, temelj, stepenice, puma, prevara, grouper, predsoblje, koza, gupi, mandarina, milisekunda, grožđe, kaiš, rog, dukserica} \\
            }
        }
    }
    & \\ \\

    \theutterance \stepcounter{utterance}  
    & & & \multicolumn{4}{p{0.6\linewidth}}{
        \cellcolor[rgb]{0.9,0.9,0.9}{
            \makecell[{{p{\linewidth}}}]{
                \texttt{\tiny{[GM$\langle$P2]}}
                \texttt{POGODAK: stepenice, puma} \\
            }
        }
    }
     \\ \\

    \theutterance \stepcounter{utterance}  
    & & & \multicolumn{2}{p{0.3\linewidth}}{
        \cellcolor[rgb]{0.9,0.9,0.9}{
            \makecell[{{p{\linewidth}}}]{
                \texttt{\tiny{[GM$|$GM]}}
                ['stepenice', 'puma']
            }
        }
    }
    & & \\ \\

    \theutterance \stepcounter{utterance}  
    & & & \multicolumn{2}{p{0.3\linewidth}}{
        \cellcolor[rgb]{0.9,0.9,0.9}{
            \makecell[{{p{\linewidth}}}]{
                \texttt{\tiny{[GM$|$GM]}}
                {'word': 'stepenice', 'assignment': 'team'}
            }
        }
    }
    & & \\ \\

    \theutterance \stepcounter{utterance}  
    & & & \multicolumn{2}{p{0.3\linewidth}}{
        \cellcolor[rgb]{0.9,0.9,0.9}{
            \makecell[{{p{\linewidth}}}]{
                \texttt{\tiny{[GM$|$GM]}}
                {'word': 'stepenice', 'assignment': 'team'}
            }
        }
    }
    & & \\ \\

    \theutterance \stepcounter{utterance}  
    & & & \multicolumn{2}{p{0.3\linewidth}}{
        \cellcolor[rgb]{0.9,0.9,0.9}{
            \makecell[{{p{\linewidth}}}]{
                \texttt{\tiny{[GM$|$GM]}}
                {'word': 'puma', 'assignment': 'team'}
            }
        }
    }
    & & \\ \\

    \theutterance \stepcounter{utterance}  
    & & & \multicolumn{2}{p{0.3\linewidth}}{
        \cellcolor[rgb]{0.9,0.9,0.9}{
            \makecell[{{p{\linewidth}}}]{
                \texttt{\tiny{[GM$|$GM]}}
                {'word': 'puma', 'assignment': 'team'}
            }
        }
    }
    & & \\ \\

    \theutterance \stepcounter{utterance}  
    & & & \multicolumn{2}{p{0.3\linewidth}}{
        \cellcolor[rgb]{0.9,0.9,0.9}{
            \makecell[{{p{\linewidth}}}]{
                \texttt{\tiny{[GM$|$GM]}}
                {'word': 'milisekunda', 'assignment': 'opponent'}
            }
        }
    }
    & & \\ \\

    \theutterance \stepcounter{utterance}  
    & & & \multicolumn{2}{p{0.3\linewidth}}{
        \cellcolor[rgb]{0.9,0.9,0.9}{
            \makecell[{{p{\linewidth}}}]{
                \texttt{\tiny{[GM$|$GM]}}
                {'hidden': {'team': ['grožđe', 'kaiš', 'temelj', 'gupi', 'mandarina'], 'innocent': ['prevara', 'dukserica', 'koza', 'predsoblje', 'grouper'], 'opponent': ['mesec', 'rog'], 'assassin': ['oposum']}, 'revealed': {'team': {'team': ['pronevera', 'viola', 'stepenice', 'puma'], 'innocent': ['kajsija', 'pikolo'], 'opponent': ['pastrmka'], 'assassin': []}, 'opponent': {'team': [], 'innocent': [], 'opponent': ['jagoda', 'tavan', 'šešir', 'magarac', 'milisekunda'], 'assassin': []}}}
            }
        }
    }
    & & \\ \\

    \theutterance \stepcounter{utterance}  
    & & \multicolumn{4}{p{0.6\linewidth}}{
        \cellcolor[rgb]{0.9,0.9,0.9}{
            \makecell[{{p{\linewidth}}}]{
                \texttt{\tiny{[P1$\langle$GM]}}
                \texttt{Reči stepenice, puma su tačno pogođene. Potez tvog saigrača se tu završio.} \\
\texttt{Sada navedite još jedan trag koji se odnosi na neke od preostalih reči vašeg tima i listu povezanih ciljnih reči. Ne zaboravite da svoj trag započnete sa 'TRAG: ', zatim pređite u novi red i započnite svoju listu ciljnih reči, odvojenih zarezima, sa 'CILJEVI: '. Napomena: neke reči su uklonjene sa lista u poređenju sa prethodnim zahtevima.} \\
\\ 
\texttt{Vaše preostale reči tima su: grožđe, kaiš, temelj, gupi, mandarina.} \\
\texttt{Preostale reči vašeg protivnika su: mesec, rog, milisekunda.} \\
\texttt{Preostale ometajuće reči su: prevara, dukserica, koza, predsoblje, grouper.} \\
\texttt{Preostale reči ubice su: oposum.} \\
            }
        }
    }
    & \\ \\

    \theutterance \stepcounter{utterance}  
    & \multicolumn{4}{p{0.6\linewidth}}{
        \cellcolor[rgb]{0.9,0.9,0.9}{
            \makecell[{{p{\linewidth}}}]{
                \texttt{\tiny{[P1$\rangle$GM]}}
                \texttt{TRAG: osnova} \\
\texttt{CILJEVI: temelj} \\
            }
        }
    }
    & & \\ \\

    \theutterance \stepcounter{utterance}  
    & & & \multicolumn{2}{p{0.3\linewidth}}{
        \cellcolor[rgb]{0.9,0.9,0.9}{
            \makecell[{{p{\linewidth}}}]{
                \texttt{\tiny{[GM$|$GM]}}
                \texttt{osnova} \\
            }
        }
    }
    & & \\ \\

    \theutterance \stepcounter{utterance}  
    & & & \multicolumn{2}{p{0.3\linewidth}}{
        \cellcolor[rgb]{0.9,0.9,0.9}{
            \makecell[{{p{\linewidth}}}]{
                \texttt{\tiny{[GM$|$GM]}}
                ['temelj']
            }
        }
    }
    & & \\ \\

    \theutterance \stepcounter{utterance}  
    & & & \multicolumn{2}{p{0.3\linewidth}}{
        \cellcolor[rgb]{0.9,0.9,0.9}{
            \makecell[{{p{\linewidth}}}]{
                \texttt{\tiny{[GM$|$GM]}}
                {'word': 'temelj', 'assignment': 'team'}
            }
        }
    }
    & & \\ \\

    \theutterance \stepcounter{utterance}  
    & & \multicolumn{4}{p{0.6\linewidth}}{
        \cellcolor[rgb]{0.9,0.9,0.9}{
            \makecell[{{p{\linewidth}}}]{
                \texttt{\tiny{[GM$\rangle$P2]}}
                \texttt{Reči stepenice, puma su tačno pogođene. Tvoj potez se tu završio.} \\
\texttt{Sada navedite još jednu listu reči razdvojenih zarezima, od najmanje 1 do najviše 1 reči, iz sledeće liste reči koje se najbolje odnose ili su najuže povezane sa rečju 'osnova'. Ne zaboravite da započnete svoj odgovor sa POGODAK:. Napomena: neke reči su uklonjene sa liste u odnosu na prethodne zahteve.} \\
\\ 
\texttt{oposum, mesec, temelj, prevara, grouper, predsoblje, koza, gupi, mandarina, grožđe, kaiš, rog, dukserica} \\
            }
        }
    }
    & \\ \\

    \theutterance \stepcounter{utterance}  
    & & & \multicolumn{4}{p{0.6\linewidth}}{
        \cellcolor[rgb]{0.9,0.9,0.9}{
            \makecell[{{p{\linewidth}}}]{
                \texttt{\tiny{[GM$\langle$P2]}}
                \texttt{POGODAK: temelj} \\
            }
        }
    }
     \\ \\

    \theutterance \stepcounter{utterance}  
    & & & \multicolumn{2}{p{0.3\linewidth}}{
        \cellcolor[rgb]{0.9,0.9,0.9}{
            \makecell[{{p{\linewidth}}}]{
                \texttt{\tiny{[GM$|$GM]}}
                ['temelj']
            }
        }
    }
    & & \\ \\

    \theutterance \stepcounter{utterance}  
    & & & \multicolumn{2}{p{0.3\linewidth}}{
        \cellcolor[rgb]{0.9,0.9,0.9}{
            \makecell[{{p{\linewidth}}}]{
                \texttt{\tiny{[GM$|$GM]}}
                {'word': 'temelj', 'assignment': 'team'}
            }
        }
    }
    & & \\ \\

    \theutterance \stepcounter{utterance}  
    & & & \multicolumn{2}{p{0.3\linewidth}}{
        \cellcolor[rgb]{0.9,0.9,0.9}{
            \makecell[{{p{\linewidth}}}]{
                \texttt{\tiny{[GM$|$GM]}}
                {'word': 'temelj', 'assignment': 'team'}
            }
        }
    }
    & & \\ \\

    \theutterance \stepcounter{utterance}  
    & & & \multicolumn{2}{p{0.3\linewidth}}{
        \cellcolor[rgb]{0.9,0.9,0.9}{
            \makecell[{{p{\linewidth}}}]{
                \texttt{\tiny{[GM$|$GM]}}
                {'word': 'mesec', 'assignment': 'opponent'}
            }
        }
    }
    & & \\ \\

    \theutterance \stepcounter{utterance}  
    & & & \multicolumn{2}{p{0.3\linewidth}}{
        \cellcolor[rgb]{0.9,0.9,0.9}{
            \makecell[{{p{\linewidth}}}]{
                \texttt{\tiny{[GM$|$GM]}}
                {'hidden': {'team': ['grožđe', 'kaiš', 'gupi', 'mandarina'], 'innocent': ['prevara', 'dukserica', 'koza', 'predsoblje', 'grouper'], 'opponent': ['rog'], 'assassin': ['oposum']}, 'revealed': {'team': {'team': ['pronevera', 'viola', 'stepenice', 'puma', 'temelj'], 'innocent': ['kajsija', 'pikolo'], 'opponent': ['pastrmka'], 'assassin': []}, 'opponent': {'team': [], 'innocent': [], 'opponent': ['jagoda', 'tavan', 'šešir', 'magarac', 'milisekunda', 'mesec'], 'assassin': []}}}
            }
        }
    }
    & & \\ \\

    \theutterance \stepcounter{utterance}  
    & & \multicolumn{4}{p{0.6\linewidth}}{
        \cellcolor[rgb]{0.9,0.9,0.9}{
            \makecell[{{p{\linewidth}}}]{
                \texttt{\tiny{[P1$\langle$GM]}}
                \texttt{Reč temelj je tačno pogođena. Potez tvog saigrača se tu završio.} \\
\texttt{Sada navedite još jedan trag koji se odnosi na neke od preostalih reči vašeg tima i listu povezanih ciljnih reči. Ne zaboravite da svoj trag započnete sa 'TRAG: ', zatim pređite u novi red i započnite svoju listu ciljnih reči, odvojenih zarezima, sa 'CILJEVI: '. Napomena: neke reči su uklonjene sa lista u poređenju sa prethodnim zahtevima.} \\
\\ 
\texttt{Vaše preostale reči tima su: grožđe, kaiš, gupi, mandarina.} \\
\texttt{Preostale reči vašeg protivnika su: mesec, rog.} \\
\texttt{Preostale ometajuće reči su: prevara, dukserica, koza, predsoblje, grouper.} \\
\texttt{Preostale reči ubice su: oposum.} \\
            }
        }
    }
    & \\ \\

    \theutterance \stepcounter{utterance}  
    & \multicolumn{4}{p{0.6\linewidth}}{
        \cellcolor[rgb]{0.9,0.9,0.9}{
            \makecell[{{p{\linewidth}}}]{
                \texttt{\tiny{[P1$\rangle$GM]}}
                \texttt{TRAG: pojas} \\
\texttt{CILJEVI: kaiš} \\
            }
        }
    }
    & & \\ \\

    \theutterance \stepcounter{utterance}  
    & & & \multicolumn{2}{p{0.3\linewidth}}{
        \cellcolor[rgb]{0.9,0.9,0.9}{
            \makecell[{{p{\linewidth}}}]{
                \texttt{\tiny{[GM$|$GM]}}
                \texttt{pojas} \\
            }
        }
    }
    & & \\ \\

    \theutterance \stepcounter{utterance}  
    & & & \multicolumn{2}{p{0.3\linewidth}}{
        \cellcolor[rgb]{0.9,0.9,0.9}{
            \makecell[{{p{\linewidth}}}]{
                \texttt{\tiny{[GM$|$GM]}}
                ['kaiš']
            }
        }
    }
    & & \\ \\

    \theutterance \stepcounter{utterance}  
    & & & \multicolumn{2}{p{0.3\linewidth}}{
        \cellcolor[rgb]{0.9,0.9,0.9}{
            \makecell[{{p{\linewidth}}}]{
                \texttt{\tiny{[GM$|$GM]}}
                {'word': 'kaiš', 'assignment': 'team'}
            }
        }
    }
    & & \\ \\

    \theutterance \stepcounter{utterance}  
    & & \multicolumn{4}{p{0.6\linewidth}}{
        \cellcolor[rgb]{0.9,0.9,0.9}{
            \makecell[{{p{\linewidth}}}]{
                \texttt{\tiny{[GM$\rangle$P2]}}
                \texttt{Reč temelj je tačno pogođena. Tvoj potez se tu završio.} \\
\texttt{Sada navedite još jednu listu reči razdvojenih zarezima, od najmanje 1 do najviše 1 reči, iz sledeće liste reči koje se najbolje odnose ili su najuže povezane sa rečju 'pojas'. Ne zaboravite da započnete svoj odgovor sa POGODAK:. Napomena: neke reči su uklonjene sa liste u odnosu na prethodne zahteve.} \\
\\ 
\texttt{oposum, prevara, grouper, predsoblje, koza, gupi, mandarina, grožđe, kaiš, rog, dukserica} \\
            }
        }
    }
    & \\ \\

    \theutterance \stepcounter{utterance}  
    & & & \multicolumn{4}{p{0.6\linewidth}}{
        \cellcolor[rgb]{0.9,0.9,0.9}{
            \makecell[{{p{\linewidth}}}]{
                \texttt{\tiny{[GM$\langle$P2]}}
                \texttt{POGODAK: kaiš} \\
            }
        }
    }
     \\ \\

    \theutterance \stepcounter{utterance}  
    & & & \multicolumn{2}{p{0.3\linewidth}}{
        \cellcolor[rgb]{0.9,0.9,0.9}{
            \makecell[{{p{\linewidth}}}]{
                \texttt{\tiny{[GM$|$GM]}}
                ['kaiš']
            }
        }
    }
    & & \\ \\

    \theutterance \stepcounter{utterance}  
    & & & \multicolumn{2}{p{0.3\linewidth}}{
        \cellcolor[rgb]{0.9,0.9,0.9}{
            \makecell[{{p{\linewidth}}}]{
                \texttt{\tiny{[GM$|$GM]}}
                {'word': 'kaiš', 'assignment': 'team'}
            }
        }
    }
    & & \\ \\

    \theutterance \stepcounter{utterance}  
    & & & \multicolumn{2}{p{0.3\linewidth}}{
        \cellcolor[rgb]{0.9,0.9,0.9}{
            \makecell[{{p{\linewidth}}}]{
                \texttt{\tiny{[GM$|$GM]}}
                {'word': 'kaiš', 'assignment': 'team'}
            }
        }
    }
    & & \\ \\

    \theutterance \stepcounter{utterance}  
    & & & \multicolumn{2}{p{0.3\linewidth}}{
        \cellcolor[rgb]{0.9,0.9,0.9}{
            \makecell[{{p{\linewidth}}}]{
                \texttt{\tiny{[GM$|$GM]}}
                {'word': 'rog', 'assignment': 'opponent'}
            }
        }
    }
    & & \\ \\

    \theutterance \stepcounter{utterance}  
    & & & \multicolumn{2}{p{0.3\linewidth}}{
        \cellcolor[rgb]{0.9,0.9,0.9}{
            \makecell[{{p{\linewidth}}}]{
                \texttt{\tiny{[GM$|$GM]}}
                {'hidden': {'team': ['grožđe', 'gupi', 'mandarina'], 'innocent': ['prevara', 'dukserica', 'koza', 'predsoblje', 'grouper'], 'opponent': [], 'assassin': ['oposum']}, 'revealed': {'team': {'team': ['pronevera', 'viola', 'stepenice', 'puma', 'temelj', 'kaiš'], 'innocent': ['kajsija', 'pikolo'], 'opponent': ['pastrmka'], 'assassin': []}, 'opponent': {'team': [], 'innocent': [], 'opponent': ['jagoda', 'tavan', 'šešir', 'magarac', 'milisekunda', 'mesec', 'rog'], 'assassin': []}}}
            }
        }
    }
    & & \\ \\

    \theutterance \stepcounter{utterance}  
    & & \multicolumn{4}{p{0.6\linewidth}}{
        \cellcolor[rgb]{0.9,0.9,0.9}{
            \makecell[{{p{\linewidth}}}]{
                \texttt{\tiny{[P1$\langle$GM]}}
                \texttt{Reč kaiš je tačno pogođena. Potez tvog saigrača se tu završio.} \\
\texttt{Sada navedite još jedan trag koji se odnosi na neke od preostalih reči vašeg tima i listu povezanih ciljnih reči. Ne zaboravite da svoj trag započnete sa 'TRAG: ', zatim pređite u novi red i započnite svoju listu ciljnih reči, odvojenih zarezima, sa 'CILJEVI: '. Napomena: neke reči su uklonjene sa lista u poređenju sa prethodnim zahtevima.} \\
\\ 
\texttt{Vaše preostale reči tima su: grožđe, gupi, mandarina.} \\
\texttt{Preostale reči vašeg protivnika su: rog.} \\
\texttt{Preostale ometajuće reči su: prevara, dukserica, koza, predsoblje, grouper.} \\
\texttt{Preostale reči ubice su: oposum.} \\
            }
        }
    }
    & \\ \\

    \theutterance \stepcounter{utterance}  
    & \multicolumn{4}{p{0.6\linewidth}}{
        \cellcolor[rgb]{0.9,0.9,0.9}{
            \makecell[{{p{\linewidth}}}]{
                \texttt{\tiny{[P1$\rangle$GM]}}
                \texttt{TRAG: citrusi} \\
\texttt{CILJEVI: mandarina, grožđe} \\
            }
        }
    }
    & & \\ \\

    \theutterance \stepcounter{utterance}  
    & & & \multicolumn{2}{p{0.3\linewidth}}{
        \cellcolor[rgb]{0.9,0.9,0.9}{
            \makecell[{{p{\linewidth}}}]{
                \texttt{\tiny{[GM$|$GM]}}
                \texttt{citrusi} \\
            }
        }
    }
    & & \\ \\

    \theutterance \stepcounter{utterance}  
    & & & \multicolumn{2}{p{0.3\linewidth}}{
        \cellcolor[rgb]{0.9,0.9,0.9}{
            \makecell[{{p{\linewidth}}}]{
                \texttt{\tiny{[GM$|$GM]}}
                ['mandarina', 'grožđe']
            }
        }
    }
    & & \\ \\

    \theutterance \stepcounter{utterance}  
    & & & \multicolumn{2}{p{0.3\linewidth}}{
        \cellcolor[rgb]{0.9,0.9,0.9}{
            \makecell[{{p{\linewidth}}}]{
                \texttt{\tiny{[GM$|$GM]}}
                {'word': 'mandarina', 'assignment': 'team'}
            }
        }
    }
    & & \\ \\

    \theutterance \stepcounter{utterance}  
    & & & \multicolumn{2}{p{0.3\linewidth}}{
        \cellcolor[rgb]{0.9,0.9,0.9}{
            \makecell[{{p{\linewidth}}}]{
                \texttt{\tiny{[GM$|$GM]}}
                {'word': 'grožđe', 'assignment': 'team'}
            }
        }
    }
    & & \\ \\

    \theutterance \stepcounter{utterance}  
    & & & \multicolumn{2}{p{0.3\linewidth}}{
        \cellcolor[rgb]{0.9,0.9,0.9}{
            \makecell[{{p{\linewidth}}}]{
                \texttt{\tiny{[GM$|$GM]}}
                \texttt{opponent has won} \\
            }
        }
    }
    & & \\ \\

\end{supertabular}
}

\end{document}
